\chapter{\textbf{INTRODUÇÃO}}

\section{\textbf{Contextualização e problema}}

A Constituição da República Federativa do Brasil de 1988 define, em seu Art. nº 194, o Sistema de Seguridade Social, com o objetivo de garantir direitos associados à saúde, à previdência e à assistência social. O pilar da previdência visa assegurar a renda ao trabalhador e seus dependentes em face de perda de força da capacidade laboral, invalidez ou morte \cite{brasil1988constituicao}.

O Regime Próprio de Previdência Social (RPPS), objeto do presente estudo, é o sistema previdenciário específico dos servidores públicos efetivos e tem suas normas básicas estabelecidas no Art. nº 40 da Constituição da República Federativa do Brasil \cite{brasil1988constituicao}. O Art. 1º da Lei nº 9.717/98 estabelece que os RPPS deverão ser organizados com base em normas gerais de contabilidade e atuária, de modo a garantir o seu equilíbrio financeiro e atuarial \cite{PL7776_2010}. Para que regimes de longo prazo, como os RPPS, possam honrar os compromissos futuros de seus planos previdenciários, é indispensável uma gestão financeira e de riscos rigorosa \cite{bogoni2011}. Em sua essência, como explica \citeonline{lopes2010impacto}, os planos previdenciários são mecanismos de proteção construídos para garantir recursos frente à perda de capacidade de trabalho.

O inciso I do Art. 1º da Lei nº 9.717/98 declara a obrigatoriedade da realização de avaliação atuarial inicial e em cada balanço utilizando-se parâmetros gerais, para a organização e revisão do plano de custeio e benefícios \cite{PL7776_2010}. Dentro da ciência atuarial, existe um processo que visa o gerenciamento e a avaliação de riscos e sistemas previdenciários, em geral. Esse processo é a avaliação atuarial. No contexto previdenciário, essa análise se empenha em assegurar a sustentabilidade dos planos por meio de precisas previsões sobre as obrigações futuras do plano, permitindo que os recursos auferidos sejam suficientes para cumprir essas obrigações futuras \cite{micoccini2010}.

Essa busca pela sustentabilidade ocorre por meio da verificação do equilíbrio do plano que, segundo \citeonline{correa2020premissas}, exige a comparação entre os pagamentos de benefícios e as contribuições em uma mesma data. É importante salientar que não se deve confundir o equilíbrio financeiro com o equilíbrio atuarial. A Portaria MTP nº 1.467/2022 distingue esses dois conceitos. O primeiro se refere à garantia de equivalência entre as receitas auferidas e as obrigações do RPPS em cada exercício financeiro. Já o segundo, se refere à garantia de equivalência, a valor presente, entre o fluxo das receitas estimadas e das obrigações projetadas, apuradas atuarialmente, a longo prazo \cite{mtp1467_2022}.

Como explica \citeonline{correa2020premissas}, as principais modalidades de planos são Benefício Definido (BD), cujo valor do benefício futuro é preestabelecido, e os de Contribuição Definida (CD), cujo benefício depende do saldo acumulado. Entretanto, ainda existe uma terceira modalidade, a Contribuição Variável (CV), que mescla características dos dois modelos principais. O fato de a maioria dos RPPS se enquadrar no modelo BD, cujo benefício é uma promessa futura, torna a precisão das premissas e do cálculo atuarial ainda mais crítica para a sustentabilidade do sistema \cite{correa2020premissas}.

Na análise de \citeonline{rodrigues2012risco}, as premissas atuariais devem representar a massa segurada da maneira mais coerente e realista possível. Portanto, quando estas não representam fidedignamente o grupo de segurados, isso pode gerar cálculos imprecisos, resultando em déficit ou superávit técnico. Conforme a classificação do mesmo, as premissas atuariais são divididas em três grupos: econômicas, biométricas e genéricas.

As premissas biométricas, conforme explica \citeonline{rodrigues2012risco}, representam os eventos como mortalidade, invalidez e sobrevivência. A escolha e a adequação dessas hipóteses são determinantes para a apuração fidedigna das responsabilidades do plano, influenciando diretamente o cálculo das provisões matemáticas e, por conseguinte, as alíquotas de contribuição necessárias para o financiamento dos benefícios \cite{correa2020premissas}.

Segundo dados do Ministério da Previdência Social referentes ao exercício de 2021, existiam 2.152 RPPS entre Estados, Distrito Federal e Municípios no Brasil. Desses, 1.732 estão em déficit atuarial, somando um total aproximado de R\$ 2,5 trilhões em déficit \cite{brasil_dados_rpps}.

A Portaria MTP Nº 1.467/2022 surge nesse contexto, estabelecendo parâmetros e diretrizes para as avaliações atuariais, com o objetivo de promover a higidez financeira desses regimes. Contudo, a aplicação prática das normativas pode suscitar debates sobre a adequação de determinados critérios.

A referida Portaria em sua Seção VI, Capítulo IV, estabelece critérios mínimos a partir do Art. 36, inciso I, e trata da verificação da validade das tábuas biométricas (usadas para estimar longevidade e entrada em invalidez) por meio da comparação entre expectativas de vida da seguinte forma:

\begin{citacao}
\noindent
[...] será dado pela tábua anual de mortalidade do Instituto Brasileiro de Geografia e Estatística -- IBGE, segregada obrigatoriamente por sexo, divulgada pela SPREV; e será averiguado por meio da comparação entre a Expectativa de Vida -- Ex estimada por essa tábua com aquela gerada pelas tábuas utilizadas na avaliação atuarial \cite[p.~28]{mtp1467_2022}.
\end{citacao}

Por meio desses itens, o Art. 36 da Portaria MTP nº 1.467/2022 estabelece que a tábua de mortalidade disponibilizada pelo IBGE será o mínimo aceitável. O RPPS poderá adotar outra tábua de mortalidade, desde que a expectativa de vida na idade média da massa não seja inferior à da tábua do IBGE.

A escolha da média como métrica central de avaliação convida a uma reflexão sobre a adequação desse critério avaliativo. Como explicam \citeonline{bruce2020estatistica}, a média artimética é uma medida de posição que, embora seja útil em um vasto leque de situações, acaba sendo uma medida frágil em casos com amostras assimétricas. Isso porque essa métrica é sensível a valores extremos e ignora a forma da distribuição (dispersão, assimetria e concentração). Portanto, a média artimética pode mascarar realidades demográficas quando usada para comparar populações.

A expectativa de vida, por sua vez, como explicam \citeonline{smith2001state}, é a duração média da vida futura a partir de uma idade específica, sendo obtida a partir da divisão entre o tempo total de vida restante de uma coorte e a número de pessoas vivas, o que a caracteriza, em termos matemáticos, como uma média aritmética. Consequentemente, ela sofre das limitações estatísticas inerentes a qualquer média.

A concentração de segurados em determinadas faixas etárias representa uma fragilidade atuarial com implicações financeiras significativas \cite{antolin2007longevity}. Em massas seguradas com distribuições etárias heterogêneas, indicadores baseados em médias podem mascarar riscos relevantes ao desconsiderar o comportamento da curva de mortalidade. Nesse contexto, uma tábua pode indicar maior sobrevida em idades anteriores ao ponto médio, mas apresentar menor sobrevida nas idades posteriores a esse ponto. Assim, caso haja uma concentração de segurados em faixas etárias nas quais a tábua apresenta menor sobrevida, ela poderá levar à subestimação da obrigação definida pelo limite mínimo, mesmo sendo, do ponto de vista matemático, superior no ponto médio. Dessa forma, duas carteiras com idade média similar podem apresentar riscos atuariais diferentes, dependendo de como os segurados se distribuem ao longo das faixas etárias \cite{rockefeller2016pension}.

Seja $PM^{i}$ a provisão matemática calculada a partir da tábua de mortalidade $i$. Esta grandeza é intrinsecamente dependente das anuidades atuariais e das probabilidades de sobrevivência associadas à referida tábua. A ideia de utilizar a expectativa de vida como indicador do nível da PM, produzida diferentes tábuas de mortalidade, se fundamenta no raciocínio de que segurados com maior longevidade esperada receberão benefícios por períodos mais extensos, elevando o valor presente das obrigações do plano \cite{silveira_santos_2017}.

Diante disso, formula-se a hipótese central desta pesquisa: o critério estabelecido pelo Art. 36 da Portaria MTP 1.467/2022 pressupõe, implicitamente, que uma tábua comparativa com expectativa de vida na idade média igual ou superior à da tábua de referência produzirá uma PM igual ou superior à da tábua de referência. Logo, sejam $e^{*}_{\overline{X}}$ a expectativa de vida na idade média segundo uma tábua comparativa qualquer e $PM^{*}$ a provisão matemática correspondente, sendo $e^{\kappa}_{\overline{X}}$ a expectativa de vida na idade média segundo a tábua de referência (IBGE 2024 extrapolada para cada gênero) e $PM^{\kappa}$ a respectiva provisão matemática. A hipótese pode ser expressa pela relação de ordem $e^{*}_{\overline{X}} \geq e^{\kappa}_{\overline{X}} \Rightarrow PM^{*} \geq PM^{\kappa}$. Contudo, ao simplificar tal hipótese usando a anuidade atuarial e a expectativa de vida na idade $x$, ou seja, avaliando $e^{*}_{x} \geq e^{\kappa}_{x} \Rightarrow a^{*}_{x} \geq a^{\kappa}_{x}$, é possível demonstrar que essa relação de ordem não é sempre verdadeira, como mostra o exemplo contido no \ref{apendice_A}.

Diante do exposto, e no contexto de massas seguradas com distribuições etárias assimétricas, o presente trabalho busca responder à seguinte questão central: \textbf{O critério de seleção de tábuas de mortalidade estabelecido pelo Art. 36 da Portaria MTP nº 1.467/2022 é suficiente para garantir que a provisão matemática calculada com a tábua adotada seja igual ou superior àquela obtida com a tábua de referência?}

\section{\textbf{Objetivos}}
Em razão desse cenário de desequilíbrio atuarial, marcado por cálculos que podem envolver distribuições assimétricas, apresentam-se os objetivos desta pesquisa.

\subsection{Objetivo geral}
Este estudo propõe-se a investigar se o critério de seleção de tábuas de mortalidade estabelecido pelo Art. 36 da Portaria MTP nº 1.467/2022 é suficiente para garantir que a provisão matemática calculada com a tábua adotada seja igual ou superior àquela obtida com a tábua de referência, identificando as condições demográficas e atuariais sob as quais podem ocorrer inconsistências e propondo um indicador de diagnóstico prévio.

\subsection{Objetivos específicos}

\begin{itemize}
    \item Simular cenários previdenciários segmentados em quatro perfis etários distintos para servir de base às projeções atuariais.

    \item Mensurar as provisões matemáticas para cada perfil estruturado, aplicando os métodos de custeio INE, CUP e PNI sob a ótica de diferentes tábuas de mortalidade.
    
    \item Identificar quais tábuas de mortalidade resultam em provisões matemáticas inferiores àquelas geradas pela tábua de referência, mesmo atendendo estritamente ao critério normativo vigente.
    
    \item Analisar as causas demográficas e financeiras que levam a essa subestimação.
    
    \item Propor um indicador analítico que permita sinalizar antecipadamente se uma determinada tábua tende a gerar uma provisão matemática inferior à da tábua de referência.
\end{itemize}
    
\section{\textbf{Justificativa}}

A Portaria MTP Nº 1.467/2022 surge em um contexto deficitário com a missão de estabelecer diretrizes a fim de promover o equilíbrio atuarial dos RPPS. A presente monografia se justifica pela necessidade de verificar a aplicabilidade do critério de validação de tábuas de mortalidade estabelecido no Art. 36. Essa justificativa se reforça diante da quantidade de RPPS que não utilizam as tábuas de mortalidade mínimas previstas nesse artigo. Segundo dados do Ministério da Previdência Social referentes ao exercício de 2025, aproximadamente 50\% dos RPPS não utilizaram a tábua mínima para a população masculina e a mesma proporção foi observada para a população feminina \cite{brasil_dados_rpps}.

Esta abordagem, embora intencionada a ser uma salvaguarda, simplifica a validação de um fenômeno complexo e distribucional como a mortalidade, especialmente em populações com distribuições etárias assimétricas. Tal critério ignora a dispersão e a forma da curva demográfica, que são cruciais para a correta apuração do passivo atuarial, cujo valor é diretamente impactado pela estrutura etária da carteira de segurados \cite{junior2019mortalidade}.

Uma tábua de mortalidade com expectativa de vida mais elevada, em teoria, implica em uma maior longevidade e por fim, um maior passivo atuarial \cite{silveira_santos_2017}. Estudos, como o de \citeonline{dias_santos_2009}, mostram que tábuas de mortalidade que projetam maior longevidade resultam em um valor de passivo atuarial significativamente superior. Se essa lógica é rompida, isto é, se uma tábua com maior expectativa de vida resultar em um passivo menor, isso indica uma falsa percepção de equilíbrio do RPPS, podendo comprometer a sustentabilidade do regime à longo prazo, onde futuros déficits que podem inviabilizar o pagamento de benefícios.

Ao quantificar o impacto financeiro desses possíveis cenários, esta monografia visa contribuir para a garantia dos direitos previdenciários dos servidores, sem onerar indevidamente as futuras gerações.

\section{\textbf{Estrutura do trabalho}}
Este trabalho está organizado em cinco capítulos. O primeiro introduz a pesquisa, estabelecendo o problema, os objetivos e a justificativa. O segundo capítulo dedica-se ao referencial teórico. O terceiro detalha a metodologia empregada na investigação. No quarto, são apresentados e discutidos os resultados obtidos. Finalmente, o quinto capítulo encerra o estudo com as considerações finais.